\noindent\textbf{Elección del procedimiento.} Con $y$ normal desde la pared y $x$ cuesta abajo, el balance es $\nu u_{yy}=-g\sin\theta$. Integrar dos veces con $u(0)=0$ y $u_y(h)=0$ da $u=(g\sin\theta/\nu)(hy-y^2/2)$. Evaluar en $h$ da la velocidad superficial e integrar de 0 a $h$ da el caudal por anchura.

\begin{enumerate}

\item El máximo del perfil está en la superficie libre sin esfuerzo.

\[\fitmath{\mathrm{us}=\frac{g\cdot s\cdot {h}^{2}}{2\cdot \nu}}\]

\[\fitmath{\mathrm{us}=\frac{(9.81)\cdot (0.5)\cdot {(0.0003)}^{2}}{2\cdot (1.007\times10^{-6})}\approx 0.2191907\,\mathrm{m/s}}\]

\item Integrar $u(y)$; no multiplicar la velocidad máxima por todo el espesor.

\[\fitmath{q=\frac{g\cdot s\cdot {h}^{3}}{3\cdot \nu}}\]

\[\fitmath{q=\frac{(9.81)\cdot (0.5)\cdot {(0.0003)}^{3}}{3\cdot (1.007\times10^{-6})}\approx 4.383813\times10^{-5}\,\mathrm{m^{2}/s}}\]

\item Aplicar el convenio de Reynolds de película pedido en el enunciado.

\[\fitmath{\mathrm{Re}=\frac{q}{\nu}}\]

\[\fitmath{\mathrm{Re}=\frac{(4.383813307\times10^{-5})}{(1.007\times10^{-6})}\approx 43.5334\,}\]

\end{enumerate}
